% !PRINTED 535
\weberpage{535}
%<<CONT
$F(x) = 0$, and when $\varphi(x) = 0$ is solved, all the roots
of $F(x)$ are thereby known. It therefore suffices if we confine
ourselves to the consideration of \gesperrt{irreducible Abelian equations}.

\quad The \gesperrt{Galois} group of an \gesperrt{Abelian} equation, whether it
be irreducible or not (provided only that its roots are distinct
from one another), has the property that in the composition
of its permutations the commutative law holds; it is
therefore a \gesperrt{commutative group}.

\quad Namely, let
\weq{4}{\alpha,\ \alpha_1 = \Theta_1(\alpha),\ \alpha_2 = \Theta_2(\alpha) \ldots \alpha_{m-1} = \Theta_{m-1}(\alpha)}
be the roots of $F(x)$, and
\weq{5}{\alpha,\ \alpha' = \Theta'(\alpha),\ \alpha'' = \Theta''(\alpha) \ldots}
the roots of the irreducible factor $\varphi(x)$ of $F(x)$
contained among them.

\quad Since, as already remarked, $\varphi(x) = 0$ is a \gesperrt{Galois} resolvent,
the group of the equation consists of the substitutions
of the field $\Omega(\alpha)$, that is, of the substitutions
\weq{6}{\sigma = (\alpha, \alpha),\ \sigma' = (\alpha, \alpha'),\ \sigma'' = (\alpha, \alpha'') \ldots}

\quad But this group obeys the commutative law; for
let
\[
\begin{aligned}
\sigma' &= (\alpha, \alpha') = [\alpha, \Theta'(\alpha)]\\
\sigma'' &= (\alpha, \alpha'') = [\alpha, \Theta''(\alpha)];
\end{aligned}
\]
then by \S.\,147
\[
\begin{aligned}
\sigma'\sigma'' &= [\alpha, \Theta'(\alpha)][\Theta'(\alpha), \Theta'\Theta''(\alpha)] = [\alpha, \Theta'\Theta''(\alpha)]\\
\sigma''\sigma' &= [\alpha, \Theta''(\alpha)][\Theta''(\alpha), \Theta''\Theta'(\alpha)] = [\alpha, \Theta''\Theta'(\alpha)].
\end{aligned}
\]
On account of (3) we therefore have $\sigma'\sigma'' = \sigma''\sigma'$, where $\sigma', \sigma''$ may be any
two of the substitutions $\sigma$. Consequently the group of the
substitutions of the field $\Omega(\alpha)$, and with it also the isomorphic
permutation group of the equation $F(x) = 0$, is \gesperrt{commutative}.

\quad The converse now also holds, but here the \gesperrt{irreducibility}
must be presupposed.

\quad \gesperrt{An irreducible equation $\varphi(x) = 0$ with commutative group is an Abelian equation.}

\quad Let $\alpha, \alpha_1, \alpha_2 \ldots \alpha_{m-1}$ be the roots of $\varphi(x) = 0$ and
$P$ the group of this equation, which is transitive because of the irreducibility of
$\varphi(x)$, and which we now moreover assume to be commutative.
Let $Q$ be the divisor of $P$ that leaves the element $0$
%CONT>>
